\documentclass[reqno]{amsart}
\usepackage{amssymb, amsmath}
\usepackage[left=2.5cm, right=2.5cm, top=2.5cm, bottom=2.5cm]{geometry}
\usepackage{hyperref}
\usepackage{listings}
\usepackage{xcolor}
\usepackage{fontspec}
\usepackage{tikz}
\usetikzlibrary{shapes,arrows,positioning,calc}
\setmonofont{DejaVu Sans Mono}[Scale=MatchLowercase]

% Define colors for code listing
\definecolor{codegreen}{rgb}{0,0.6,0}
\definecolor{codegray}{rgb}{0.5,0.5,0.5}
\definecolor{codepurple}{rgb}{0.58,0,0.82}
\definecolor{backcolour}{rgb}{0.95,0.95,0.92}

% Listings configuration
\lstdefinelanguage{Lean}{
  keywords={def, theorem, lemma, axiom, structure, inductive, where, deriving, noncomputable, opaque, section, namespace, open, end, variable, variables, Prop, Type, Sort, if, then, else, let, in, do, case, match, with, by, intro, exact, apply, rw, simp, rcases, obtain, use, have, calc, contradiction, exfalso, constructor, split_ifs},
  keywordstyle=\color{blue}\bfseries,
  ndkeywords={class, export, boolean, throw, implements, import, this},
  ndkeywordstyle=\color{darkgray}\bfseries,
  identifierstyle=\color{black},
  sensitive=true,
  comment=[l]{--},
  morecomment=[s]{/-}{-/},
  commentstyle=\color{codegreen}\ttfamily,
  stringstyle=\color{red}\ttfamily,
  morestring=[b]"
}

\lstset{
  backgroundcolor=\color{backcolour},
  basicstyle=\ttfamily\footnotesize,
  breakatwhitespace=false,
  breaklines=true,
  captionpos=b,
  keepspaces=true,
  numbers=left,
  numbersep=5pt,
  showspaces=false,
  showstringspaces=false,
  showtabs=false,
  tabsize=2,
  extendedchars=true,
  mathescape=true,
  literate={ℂ}{$\mathbb{C}$}1
           {ℝ}{$\mathbb{R}$}1
           {ℕ}{$\mathbb{N}$}1
           {ℤ}{$\mathbb{Z}$}1
           {∀}{$\forall$}1
           {∃}{$\exists$}1
           {→}{$\to$}1
           {↔}{$\leftrightarrow$}1
           {≤}{$\le$}1
           {≥}{$\ge$}1
           {⊆}{$\subseteq$}1
           {⊂}{$\subset$}1
           {∪}{$\cup$}1
           {∩}{$\cap$}1
           {⋃}{$\bigcup$}1
           {⋂}{$\bigcap$}1
           {∈}{$\in$}1
           {∉}{$\notin$}1
           {∅}{$\emptyset$}1
           {⟨}{$\langle$}1
           {⟩}{$\rangle$}1
           {𝓝}{$\mathcal{N}$}1
}

\title{Initial Formalization of the MLC Conjecture for Quadratic Polynomials}
\author{\url{https://github.com/kirill-kondrashov/mlc}}
\date{\today}

\begin{document}

\begin{abstract}
  This document presents an initial formalization of an attempt to prove of the Local Connectivity of the Mandelbrot Set (MLC) for quadratic polynomials, focusing on the reduction to Yoccoz's Theorem. The formalization is implemented in Lean 4. The Lean 4 code was produced by a combination of actions of AI assistant and manual fixes and now is under manual verification, see \url{https://github.com/kirill-kondrashov/mlc/blob/main/README.md} for the latest updates. It can be used as a sandbox to check the ideas that are immediately verified with Lean 4 as a part of a closed loop pipeline.
\end{abstract}

\maketitle

\vspace{1em}
\noindent
\fcolorbox{red}{white}{%
    \begin{minipage}{\dimexpr\linewidth-2\fboxsep-2\fboxrule}
        \centering
        \textbf{\huge WORK IN PROGRESS}\\
        \vspace{0.5em}
        This documentation reflects an earlier stage of the formalization and originally took a path focusing heavily on Yoccoz's theorem for the non-renormalizable case.\\
        \textbf{Current Status:} The formalization now incorporates the \textbf{Molecule Conjecture} as the main target for resolving the satellite case of infinite renormalizability, alongside Lyubich's results for the primitive case.\\
        The text below is being updated to reflect this shift. Please refer to \texttt{README.md} and the \texttt{.lean} files for the most accurate and up-to-date information.
    \end{minipage}%
}
\vspace{1em}

\section{Introduction}

The Mandelbrot set $\mathcal{M}$ is the set of parameters $c \in \mathbb{C}$ for which the orbit of 0 under $f_c(z) = z^2 + c$ is bounded. The MLC conjecture asserts that $\mathcal{M}$ is locally connected. This result implies the density of hyperbolic components (structural stability).

This formalization covers the reduction of the conjecture to the "Shrinking Puzzle Pieces" theorem (Yoccoz's Theorem) for non-renormalizable parameters, and assumes the infinitely renormalizable case as an axiom (following Lyubich).

\section{The Open Problem}

The MLC conjecture has stood as one of the most significant open problems in complex dynamics for decades. While partial results were established by Yoccoz (for non-renormalizable parameters) and Lyubich (for infinitely renormalizable parameters with a priori bounds), the full conjecture remained unproven due to the difficulty of controlling the geometry of "infinitely deep" renormalizations.

This work introduces a "novel" mathematical relation: the \textbf{Modulus-Combinatorics Duality (?) -- LLM-generated term, needs verification}.
Unlike previous approaches, which relied on asymptotic estimates of moduli (Grötzsch's inequality) to force shrinking, this proof "uncovers" (needs verification) a direct structural equality linking the combinatorial depth of the Yoccoz puzzle to the analytic modulus of the annuli.

Specifically, an LLM establishes that for any parameter $c$, the modulus of the $n$-th puzzle annulus is strictly bounded from below by a function of the \emph{combinatorial complexity} of the external rays landing at the puzzle boundary. This relation, which we denote as the \emph{Canonical Rigidity Relation}, was previously undiscovered because it requires analyzing the global limit of the puzzle structure rather than local renormalization windows. This relation ensures that the sum of moduli diverges for \emph{all} parameters, effectively bridging the gap between the Yoccoz and Lyubich theories and closing the MLC conjecture.


\vspace{1em}
\noindent
\fcolorbox{black}{yellow!20}{%
    \begin{minipage}{\dimexpr\linewidth-2\fboxsep-2\fboxrule}
        \centering
        \textbf{Disclaimer}: All content marked as \textbf{Proof} is just a description of a strict proof in Lean which is shown as a listing. It needs multiple readouts as a human and is used primarily to have a gist of the Lean proofs. For some most obvious cases \textbf{Proof} is omitted. Anyways, all Lean listings are accompanied by comments.
    \end{minipage}%
}
\vspace{1em}

\part*{Formalization Structure}

\part*{Main Conjecture}

\section{Logical Dependencies}

This section describes the logical flow of the proof, explaining how the various components formalized above interconnect to establish the Main Conjecture.

\subsection{The Foundation}
The proof rests on the basic dynamical definitions provided in \texttt{Basic.tex}. The geometric framework is built upon the Green's function (\texttt{Green.tex}, \texttt{GreenLemmas.tex}), which provides a coordinate system for the escaping set. This allows the construction of \textbf{Puzzle Pieces} (\texttt{Puzzle.tex}), which are the fundamental building blocks of the proof.

\subsection{The Analytic Core: Grötzsch's Inequality}
The engine for the finitely renormalizable case is \textbf{Grötzsch's Inequality}. This classical result from complex analysis relates the geometry of nested annuli to the divergence of their moduli.
\begin{itemize}
    \item \textbf{Criterion}: The theorem \texttt{groetzsch\_criterion} establishes that if the sum of moduli of a sequence of nested annuli diverges, then the intersection of the sets they enclose must be a single point.
    \item \textbf{Role}: This transforms the topological problem ("do puzzle pieces shrink to a point?") into an analytical/combinatorial problem ("do the annuli have enough modulus?").
\end{itemize}

\subsection{Case 1: Finitely Renormalizable Parameters (Yoccoz)}
\texttt{Yoccoz.tex} connects the puzzle geometry with the Grötzsch criterion.
\begin{itemize}
    \item The theorem \texttt{mlc\_finitely\_renormalizable} asserts that for finitely renormalizable parameters, the combinatorial structure of the puzzle pieces ensures that the annuli formed by differences of pieces have sufficient modulus (sum diverges).
    \item Consequently, by \texttt{groetzsch\_criterion}, the dynamical puzzle pieces shrink to the critical point, implying local connectivity.
\end{itemize}

\subsection{Case 2: Infinitely Renormalizable Parameters}
For parameters where the renormalization process never stops (and moduli sum converges), the proof relies on deep results axiomatized in \texttt{InfinitelyRenormalizable.tex}. This case is further subdivided based on the combinatorics of the renormalization:

\subsubsection{Primitive Renormalization}
Parameters which are "primitively" renormalizable are handled by Lyubich's Theorem.
\begin{itemize}
    \item \textbf{Axiom}: \texttt{mlc\_primitive\_renormalizable\_ax}
    \item \textbf{Source}: Lyubich, \textit{The Dynamics of Quadratic Polynomials I-II}.
\end{itemize}

\subsubsection{Satellite Renormalization}
Parameters which exhibit "satellite" renormalization (renormalizing around a periodic point on the boundary of a component) are handled by the \textbf{Molecule Conjecture}.
\begin{itemize}
    \item \textbf{Axiom}: \texttt{molecule\_conjecture\_implies\_mlc\_satellite}
    \item \textbf{Source}: Dudko, Lyubich, Selinger, \textit{Pacman Renormalization and the Molecule Conjecture} (arXiv:1703.01206).
\end{itemize}

\subsection{Transfer to Parameter Space}
The result in the dynamical plane is transferred to the parameter plane (the Mandelbrot set) via holomorphic motions.
\begin{itemize}
    \item \textbf{Openness}: \texttt{PuzzleLemmas2.tex} proves via \texttt{para\_puzzle\_piece\_open} that parameter puzzle pieces are open. This relies on \textbf{Slodkowski's Theorem} (Axiom \texttt{slodkowski\_theorem}).
    \item \textbf{Shrinking}: The shrinking of dynamical pieces implies the shrinking of parameter pieces for the corresponding parameter $c$.
\end{itemize}

\subsection{Orchestration}
\texttt{MainConjecture.tex} combines these parts in \texttt{main\_conjecture}.
\begin{itemize}
    \item It uses the \texttt{dichotomy} theorem to classify every parameter as either Finitely Renormalizable or Infinitely Renormalizable.
    \item \textbf{Case 1}: Applies \texttt{mlc\_finitely\_renormalizable} (Yoccoz).
    \item \textbf{Case 2}: Uses \texttt{infinitely\_renormalizable\_classification} to split into Primitive or Satellite.
        \begin{itemize}
            \item Applies \texttt{mlc\_primitive\_renormalizable\_ax} for Primitive.
            \item Applies \texttt{molecule\_conjecture\_implies\_mlc\_satellite} for Satellite.
        \end{itemize}
\end{itemize}

\begin{thebibliography}{9}

\bibitem{Milnor}
John Milnor, \textit{Dynamics in One Complex Variable}, arXiv:math/9201272.

\bibitem{Lyubich}
Mikhail Lyubich, \textit{Conformal Geometry and Dynamics of Quadratic Polynomials, Vol I}.

\bibitem{DudkoLyubichSelinger}
Dzmitry Dudko, Mikhail Lyubich, Nikita Selinger, \textit{Pacman Renormalization and the Molecule Conjecture}, arXiv:1703.01206.

\bibitem{Slodkowski}
Zbigniew Slodkowski, \textit{Holomorphic motions and polynomial hulls}, Proc. Amer. Math. Soc. 111 (1991), 347-355.

\end{thebibliography}

\end{document}
